\section{Conclusion and Future Work}
\label{sec:conclusion}

In this paper, we introduced \texttt{Hetero-KVCache-Optimizer}, a hierarchical KV Cache management system that enables long-context MLLM inference on resource-constrained edge devices. Our key innovation lies in the \textit{transient interception} architecture: decoupling the logical computation view from the physical memory layout to satisfy FlashAttention's structural constraints while achieving sub-linear VRAM growth.

Our three-tier storage hierarchy (HBM resident Sink+Tail, DRAM-compressed overflow, asynchronous prefetching) fundamentally differs from eviction-based approaches like StreamingLLM. By compressing rather than discarding intermediate tokens, we preserve the ability to reconstruct any portion of the sequence, achieving perfect recall on long-range dependency tasks.

\textbf{Summary of Results.} Evaluated on Qwen2-VL-7B-Instruct (4-bit NF4) with multimodal inputs up to 12K tokens:
\begin{itemize}
    \item Peak VRAM constrained to \textbf{2.65GB} at 12K multimodal tokens where native encounters OOM beyond 8K
    \item Needle-in-a-Haystack recall: \textbf{100\%} (Hetero-KV) with multimodal contexts
    \item Successfully processes \textbf{12K+ multimodal tokens} within 16GB VRAM constraint where native baseline fails with OOM
\end{itemize}

\textbf{Limitations.} We identify the following limitations of our current work:

\begin{enumerate}
    \item \textbf{Hardware Bandwidth Gap:} While we constrain the VRAM to 16GB, the A100 GPU features PCIe Gen4 bandwidth ($\sim$64GB/s) that far exceed typical edge devices (e.g., RTX 4090 with PCIe Gen4 x16). The prefetching overhead (TPOT) may increase on true edge hardware due to lower PCIe bandwidth.
    
    \item \textbf{Single-Device Deployment:} Current implementation targets single-device deployment. Extending the hierarchical storage to span multiple GPUs and edge-cloud boundaries would enable even longer context support.
    
    \item \textbf{Speculative Decoding Integration:} Combining Hetero-KV with speculative decoding could further reduce TTFT while maintaining memory efficiency.
    
    \item \textbf{Adaptive Quantization:} Implementing mixed-precision quantization (FP8/INT4/INT2) based on token importance could push context limits toward 100K+ tokens on 16GB hardware.
\end{enumerate}

We believe the ``Transient-to-Static'' architectural pattern presented in this work has broader applicability beyond KV Cache management, potentially informing future designs for activation checkpointing and gradient compression in resource-constrained training scenarios.
